\begin{helper}[Product Model Pushforward Identity]
Fix \(n\in\mathbb N\), \(I\subseteq\mathrm{Fin}(n)\), parameters \(\varepsilon\in\mathbb R\).
Let \(M\in\mathbb N\) and \(p_0\in\mathbb R\) with \(0\le p_0\le1\).
Let \(\pi:\mathrm{Fin}(n)\hookrightarrow\mathrm{Fin}(M)\times\mathrm{Fin}(M)\) be an embedding.
Let
\[
P_R=\{\pi(u)_1:u\in I\},\qquad P_B=\{\pi(u)_2:u\in I\},
\]
\[
q=\noderef{FixedSetProductModelMassFn}(p_0,P_R,P_B),
\qquad
Z_{\mathrm{prod}}=\noderef{FixedSetProductModelVar}(I,\varepsilon,p_0,\pi).
\]
Let \(Z_I\) be the sample-based quantity defined identically as in \(\noderef{FixedSetProductModelTail}\).
Then for any predicate \(H\), the conditional distribution on the cell \(\mathcal C_\pi = \{\omega:\noderef{TwoBiteEmbedding}(\omega)=\pi\}\) pushes forward to the product distribution:
\[
\begin{aligned}
&\noderef{TwoBiteEventProbability}(n,M,p_0)
 \bigl(\mathcal C_\pi\cap \{\omega:H(Z_I(\omega))\}\bigr)\\
&\quad=
\noderef{TwoBiteEventProbability}(n,M,p_0)(\mathcal C_\pi)
 \sum_{\{v:H(Z_{\mathrm{prod}}(v))\}}\prod_iq(v_i).
\end{aligned}
\]
\end{helper}
\begin{proof}
By the hypotheses \(0\le p_0\le1\),
\(\noderef{TwoBiteSampleWeightNonnegative}\) gives nonnegativity of every
two-bite sample weight, and \(\noderef{FixedSetProductModelMass}\) gives
\[
q(a)\ge0\quad\text{for every }a,\qquad \sum_a q(a)=1.
\]
Thus all sample sums and all product-assignment sums below are finite sums of
nonnegative weights.

Let \(\mathcal C_\pi=\{\omega:\noderef{TwoBiteEmbedding}(\omega)=\pi\}\).
If \(\mathcal C_\pi\) has zero probability, then
\(\mathcal C_\pi\cap\{\omega:H(Z_I(\omega))\}\subseteq\mathcal C_\pi\), so
the left side is \(0\) by nonnegativity.  The right side is the same zero cell
mass multiplied by a nonnegative product sum, hence is also \(0\).

Assume now that the cell mass is positive.  On \(\mathcal C_\pi\), the
embedding is fixed, so for every \(u\in I\) the red and blue projections are
\(\pi(u)_1\in P_R\) and \(\pi(u)_2\in P_B\).  Unfolding
\(\noderef{TwoBiteEventProbability}\), a sample in the cell has weight
\[
\noderef{GnpGraphWeight}(p_0,G_R)\,
\noderef{GnpGraphWeight}(p_0,G_B)\,
\noderef{UniformInjectionWeight}(n,M),
\]
and the injection factor is constant across the cell.

Partition the red graph edge variables into three classes:
\[
\{r_1,r_2\}\subseteq P_R,\qquad
\{x_1,x_2\}\subseteq \mathrm{Fin}(M)\setminus P_R,\qquad
\{x,r\}\text{ with }x\notin P_R,\ r\in P_R.
\]
The first two classes are ignored by \(Z_I\), because
\(\noderef{TwoBiteX}(\omega,I,\operatorname{inl}x)\) tests adjacency only
between \(x\) and red projections of vertices in \(I\), and those projections
all lie in \(P_R\).  Summing over each ignored red edge gives the factor
\[
(1-p_0)+p_0=1.
\]
The surviving red incident data are the subsets
\[
A_R(x)=\{r\in P_R:\{x,r\}\text{ is a red edge}\}
\qquad(x\notin P_R),
\]
with total red weight
\[
\prod_{x\notin P_R}
p_0^{|A_R(x)|}(1-p_0)^{|P_R|-|A_R(x)|}.
\]
The blue graph has the identical partition with \(P_B\).  Its surviving data
are
\[
A_B(y)=\{b\in P_B:\{y,b\}\text{ is a blue edge}\}
\qquad(y\notin P_B),
\]
with weight
\[
\prod_{y\notin P_B}
p_0^{|A_B(y)|}(1-p_0)^{|P_B|-|A_B(y)|}.
\]

Let \(\mathcal D\) be the finite set of all such incident-data families
\[
D=\bigl((A_R(x))_{x\notin P_R},(A_B(y))_{y\notin P_B}\bigr).
\]
Grouping samples in \(\mathcal C_\pi\) by \(D\), and summing out the ignored
red and blue graph edges, gives the cell mass contribution
\[
\noderef{UniformInjectionWeight}(n,M)
\prod_{x\notin P_R}
p_0^{|A_R(x)|}(1-p_0)^{|P_R|-|A_R(x)|}
\prod_{y\notin P_B}
p_0^{|A_B(y)|}(1-p_0)^{|P_B|-|A_B(y)|}
\]
for the data \(D\).

Now group product assignments by the same data.  A product assignment has
\(2M\) coordinates, each a pair \((R,B)\).  For a red coordinate
\(x\notin P_R\), set \(R=A_R(x)\); the second component \(B\) is redundant.
For this fixed \(R\subseteq P_R\),
\[
\sum_{B\subseteq P_B}q(R,B)
=p_0^{|R|}(1-p_0)^{|P_R|-|R|},
\]
using the definition of \(q\), the support condition in
\(\noderef{FixedSetProductModelMassFn}\), and the normalization supplied by
\(\noderef{FixedSetProductModelMass}\) in the \(P_B\)-component.  If
\(x\in P_R\), then \(\noderef{FixedSetFakeGraph}\) does not read the red
coordinate indexed by \(x\), and the whole pair \((R,B)\) sums to \(1\) by
\(\noderef{FixedSetProductModelMass}\).  Symmetrically, for a blue coordinate
\(y\notin P_B\) the second component is fixed as \(B=A_B(y)\), the first
component sums out, and the contribution is
\[
p_0^{|B|}(1-p_0)^{|P_B|-|B|};
\]
if \(y\in P_B\), the whole blue coordinate has total \(q\)-mass \(1\).

Thus the finite map from product assignments to \(\mathcal D\), obtained by
remembering only first components at red outside coordinates and second
components at blue outside coordinates, has fiber mass exactly equal to the
red and blue incident-data weight displayed above, without the constant
injection factor.  This is the required finite reindexing: both the cell
sample sum and the product-assignment sum are sums over the same finite data
set \(\mathcal D\), with the same data weight after the cell's constant
injection factor is pulled out.

These data weights have total mass \(1\).  On the graph side, summing over
all possible incident subsets \(A_R(x)\subseteq P_R\) for a fixed outside red
vertex \(x\) gives
\[
\sum_{A\subseteq P_R}p_0^{|A|}(1-p_0)^{|P_R|-|A|}
=(p_0+1-p_0)^{|P_R|}=1,
\]
and similarly for every outside blue vertex.  On the product side, the same
normalization is exactly the coordinate normalization from
\(\noderef{FixedSetProductModelMass}\).  Hence the mass of the full cell
\(\mathcal C_\pi\) is the fixed injection factor times \(1\), and the
contribution of a data value \(D\) to the cell is
\[
\noderef{TwoBiteEventProbability}(n,M,p_0)(\mathcal C_\pi)\,\nu(D),
\]
where \(\nu(D)\) is the corresponding product fiber mass.

For data \(D\), the sample graph and the product fake graph have the same
adjacencies from outside \(P_R\) into \(P_R\) and from outside \(P_B\) into
\(P_B\).  The definitions of \(\noderef{TwoBiteSmallClass}\),
\(\noderef{TwoBiteX}\), \(\noderef{TwoBiteFinalPairs}\), and
\(\noderef{FixedSetProductModelVar}\) use only these adjacencies and the fixed
embedding \(\pi\).  Hence all samples and product assignments with the same
data \(D\) have a common value \(Z(D)\), with
\[
Z_I(\omega)=Z(D)=Z_{\mathrm{prod}}(v).
\]
Grouping both sides of the desired identity by \(D\) gives
\[
\sum_{D\in\mathcal D,\ H(Z(D))}
\noderef{TwoBiteEventProbability}(n,M,p_0)(\mathcal C_\pi)\,\nu(D)
=
\noderef{TwoBiteEventProbability}(n,M,p_0)(\mathcal C_\pi)
\sum_{D\in\mathcal D,\ H(Z(D))}\nu(D),
\]
where \(\nu(D)\) is the product fiber mass of \(D\).  This is exactly
\[
\noderef{TwoBiteEventProbability}(n,M,p_0)
 \bigl(\mathcal C_\pi\cap \{\omega:H(Z_I(\omega))\}\bigr)
=
\noderef{TwoBiteEventProbability}(n,M,p_0)(\mathcal C_\pi)
 \sum_{\{v:H(Z_{\mathrm{prod}}(v))\}}\prod_iq(v_i),
\]
which proves the pushforward identity.
\end{proof}
